\begin{frame}
\frametitle{What is resource in undertype?}

We are actually verifying the following VCs: 
\begin{align*}
\forall x_1. \exists x_2. \forall x_3. \exists x_4. e : \tau 
\end{align*}

We interpret $\forall x$ as ``read permission'', i.e., someone decides the value of 
$x$ and we can only read it. We interpret $\exists x$ as ``write permission'', i.e., we decide the value of $x$ that we want.

Why different?

\end{frame}

\begin{frame}
    \frametitle{Resource Interpretation}
    
Consider the following example:
\begin{align*}
    \Gamma \vdash x - y : \rt{\Int}{\vnu = 0}
\end{align*}
\begin{itemize}
    \item If we have write permission for both $x$ and $y$, we can set $x = 3$ and $y = 3$, then the type judgment is valid.
    \item If we have read permission for $x$ and write permission for $y$, we can set $y = \Code{read}(x)$, then the type judgment is valid.
    \item If we only have read permission for both $x$ and $y$, we cannot do anything. Since we cannot control the value we read, we cannot prove the type judgment.
    \item If we have write permission for $x$ and read permission for $y$, we still cannot do anything. Since the read happens \textbf{after} the write, we cannot ``guess'' the value of $y$.
\end{itemize}
    
\end{frame}

\begin{frame}
    \frametitle{Terminology}
    
The bunched type and the bunched implication share the same key idea (combine two modalities within one system) but one in type theory and one in logic. 

Bunched type plays lots of games with type context (it has two unit contexts for two different monoids), which makes it much less intuitive than separation logic. 

I also assume you are more familiar with separation logic, so I will use separation logic to explain my idea.

Note that classical separation logic (bunched implication, bunched type) is incompariable with linear type, thus I won't use their terminology (e.g., number of resources one can use). Although introducing the $!$ operator is reasonable (persistently modality in Iris), I don't do that for now.
    
\end{frame}

\begin{frame}
    \frametitle{Terminology: Variables and Heaps}
    
In SL, we have two worlds: pure logic world ($x = 4$), and heap world ($x \mapsto 4$). One can be duplicated and the other cannot. We also introduce two forms of type annotations:
\begin{itemize}
    \item $x : \tau$ similar to $x = 4$: we have read permission to get an assignment of $x$ consistent with $\tau$.
    \item $x \mapsto \tau$ similar to $x \mapsto 4$: we have write permission to decide the assignment of $x$ to a value consistent with $\tau$. (For now we assume $\mapsto$ only works with base type)
\end{itemize}

Similar to SL, we can duplicate $x : \tau$ (I can read multiple times) but not $x \mapsto \tau$ (I can only decide once).

Note that this permission is not ``affine'', which means we cannot give up the write permission. If we want to type check a function $\lambda x.e$ whose under-style parameter type is $\rt{b}{\bot}$, we cannot find such an assignment for $x$ (and we cannot give up), thus the type check should fail.
\end{frame}

\begin{frame}
    \frametitle{Terminology: Actions and Randomness}

Each action means one truly random choice.

The right way to precisely define it is that the player action is a deterministic function (winning strategy) 
\[ f: \overline{x} \to \Code{seed} \to \overline{y} \]
where $\overline{x}$ are the variables to read and $\overline{y}$ are the variables to write. The result of written values depends on the random seed.

\end{frame}

\begin{frame}
    \frametitle{Terminology: Connectives}
    
In SL, we have two connectives: $\land$ and $\sep$. 

$P \sep Q$ means heap $P$ and $Q$ are disjoint, however, $\land$ doesn't guarantee this disjointness. In our two-player game, one player can make player actions (i.e., decide a value to write); heap means action, and disjoint means in two different actions.

We also introduce two similar connectives in UnderType: $,$ and $\sep$.

$\Delta_1, \Delta_2$ is similar to $P \land Q$, it means that $\Delta_1$ and $\Delta_2$ are in the same player action. Note that $\Delta_1, \Delta_2$ is commutative only when $\Delta_2$ does not refer to the variables in $\Delta_1$.

$\Gamma_1 \sep \Gamma_2$ is similar to $P \sep Q$, it means that $\Gamma_1$ and $\Gamma_2$ are in different player actions.
\end{frame}

\begin{frame}
    \frametitle{Terminology: Connectives}
    
Following this idea, we normalize the type context language into a more restricted form:
\begin{align*}
    &\text{heap} \quad \Delta \doteq \overline{x{:}\tau},\Code{emp} ~|~ \overline{x{:}\tau}, y\mapsto \tau
    \\&\text{type context} \quad \Gamma \doteq \Delta ~|~ \Gamma \sep \Gamma
\end{align*}

Each (nonempty) heap (action) means reading some variables ($\overline{x}$) that are connected with $,$ and then writing to a variable ($y$). 

We also have empty heap $\Code{emp}$ which means no write operations; it is for a pure over-approximated type context.

\end{frame}

\begin{frame}
    \frametitle{Terminology: Connectives}
    
Intuitively, this definition forces dividing each potential $\mapsto$ operation into different heaps. This disallows the following case:
\begin{align*}
    y_1 \mapsto \Int, y_2 \mapsto \Int
\end{align*}

The meaning of this type context is: we want to write both $y_1$ and $y_2$, however they probably are in the same action. According to the definition of action, it means that $y_1$ and $y_2$ may depend on the same random seed, thus they are not truly random (independent).

\end{frame}

\begin{frame}
    \frametitle{Terminology: Connectives}
    
Compared with SL, $x \mapsto 3 \land y \mapsto 4$ in standard SL is impossible since $x \mapsto 3$ means this is a single address in the current heap. The $\land$ means both $x$ and $y$ are the same address and $3 = 4$.

In SL, we hardly see a formula like this.

In our setting $x \mapsto \Int \land y \mapsto \Int$, we don't have a universal namespace (store) and treat variable names more flexibly, i.e., $x$ and $y$ can be different variables with different assignments. It means they are entangled in an unknown way (e.g., there exists an unknown deterministic function $g$ where $g(x) = y$).

\end{frame}

\begin{frame}
    \frametitle{Terminology: Connectives}

Although our interpretation allows $x \mapsto \Int \land y \mapsto \Int$, it is a little bit weird to write two variables in the same action, since in our incorrectness reasoning setting, all parameters can have independent random choices. One workaround is 
\begin{align*}
    y_1: \Int, y_2 \mapsto \Int
\end{align*}
where we make $y_1$ fall back to a read variable.

\end{frame}


\begin{frame}
    \frametitle{Structural Rules: Exchange}   

\begin{figure}[!ht]
    \mprooftr{30mm}{
     $\Delta_2, \Delta_1 : e : \tau$ \quad
     $\wfvdash \Delta_1$ 
    }{ExchangeO}{
     $\Delta_1, \Delta_2 : e : \tau$}
\quad
\mprooftr{30mm}{
     $\Gamma_2 \sep \Gamma_1 : e : \tau$
    }{ExchangeU}{
     $\Gamma_1 \sep \Gamma_2 : e : \tau$}
\end{figure}

The $,$ connective is commutative only when $\Delta_1$ doesn't refer to the variables in $\Delta_2$.
    
\end{frame}

\begin{frame}
    \frametitle{Structural Rules: Contraction}   

\begin{figure}[!ht]
    \mprooftr{40mm}{
     $\Delta_1, \Delta_1, \Delta_2 \vdash : e : \tau$ 
    }{ContractionLeft}{
     $\Delta_1, \Delta_2 : e : \tau$}
\quad
\mprooftr{40mm}{
    $\Delta_1, \Delta_2, \Delta_2 \vdash : e : \tau$ \quad
    $\mapsto \notin \Delta_2$
    }{ContractionRight}{
     $\Delta_1, \Delta_2 : e : \tau$}
\quad
\mprooftr{40mm}{
    $\Gamma \sep (\Code{Over}(\Gamma), \Code{emp})\vdash : e : \tau$ 
    }{ContractionEmp}{
     $\Gamma : e : \tau$}
\end{figure}

The contraction rule defines what we can duplicate; any $x \mapsto \tau$ can be duplicated. The $\Code{emp}$ and read variables can be duplicated. 
Here $\Code{Over}(\Gamma)$ means we keep all variables in $\Gamma$ with read permission.
\end{frame}

\begin{frame}
    \frametitle{Structural Rules: Weakening}   

\begin{figure}[!ht]
    \mprooftr{40mm}{
     $\Delta_2 \vdash : e : \tau$ 
    }{WeakeningLeft}{
     $\Delta_1, \Delta_2 : e : \tau$}
\quad
\mprooftr{40mm}{
    $\Delta_1 \vdash : e : \tau$ \quad
    $\mapsto \notin \Delta_2$
    }{WeakeningRight}{
     $\Delta_1, \Delta_2 : e : \tau$}
\quad
\mprooftr{50mm}{
    $\Gamma \vdash : e : \tau$ \quad
    $\Gamma'$ is not empty
    }{Frame}{
     $\Gamma \sep \Gamma' : e : \tau$}
\end{figure}

The weakening rule defines what we can drop, which is similar to contraction rules. One difference is that the weakening rule for $\sep$ (frame rule) requires that the dropped context is not empty (e.g., $x\mapsto\rt{\Int}{\bot}$).
\end{frame}

\begin{frame}[fragile]
    \frametitle{Typing}   

I would like to solve the following typing examples, where $f$ and $g$ can be assigned with different function types:

\begin{minted}{ocaml}
let y1 = f x in
let y2 = g x in
x
\end{minted}

We want to show:
\begin{itemize}
    \item if $f$ and $g$ both have under-parameters, this typecheck should fail.
    \item if $f$ and $g$ both have over-parameters, this typecheck should succeed.
    \item if $f$ and $g$ one has over-parameter and the other has under-parameter, this typecheck should succeed. 
    It looks like we separate a read permission to the over-parameter function, and later recover the write permission back.
\end{itemize}

\end{frame}

\begin{frame}[fragile]
    \frametitle{Function Typing}   

Before talking about the typing, we introduce the function types.

Similar to how we have normal implication ($\impl$) and magic wand ($\wand$, it is actually linear arrow) in separation logic, we have two arrow types:
\begin{align*}
    &\tau_1 \doteq x{:}\Int \sarr \Int \quad (over-to-under, coverage type)
    \\&\tau_2 \doteq (x\mapsto\Int) \wand \Int \quad (under-to- under, incorrectness triple)
\end{align*}

How should we interpret these two types? 

The first means if we read arbitrary integer value for $x$, we can get a random integer value (or a new permission to make a decision). 

The second means permission transfer, we consume the permission to decide $x$, but get a fresh permission.

\end{frame}

\begin{frame}[fragile]
    \frametitle{Other potential function types}   

There are two types allowed by the syntax
\begin{align*}
    &\tau_1 \doteq x{:}\Int \wand \Int 
    \\&\tau_2 \doteq (x\mapsto\Int) \sarr \Int 
\end{align*}

The first type requires the ``separation'', but its argument is pure (based on the empty action $\Code{emp}$), thus it is equal to $x{:}\Int \sarr \Int$. (Similar in SL).

The second type requires no ``separation'', however connected two actions, which may lead the type context like $x\mapsto\Int, y\mapsto\Int$ -- we already shown some examples about it, it is quite annoying thus I disallow it for now.

\textbf{Remark:} I simplify all parameters in form $x{:}\tau_x$ from now, where the relation between $x$ and $\tau_x$ is clear according to the later types. (Actually, the bunched type can only use $,$ and $\sep$ to distinguish all $x:\tau$ and $x\mapsto\tau$ even in type context, but it is not clear for presentation).

\end{frame}

\begin{frame}[fragile]
    \frametitle{Missing function types}   

The ``over'' returned function type is missing, it can be expressed as:
\begin{align*}
    &f_1: x{:}\urt{\Int}{\top} \sarr \ort{\Int}{\top} \doteq x{:}\Int \wand (\Int, \Code{emp})
    \\&f_2: x{:}\urt{\Int}{\top} \sarr \urt{\Int}{\top} \doteq x{:}\Int \wand \Int
\end{align*}

Here we highlight the return value ($\Int$) has only read permission (only pure constraint can connected with $\Code{emp}$ with $,$). To show the difference we (informally) have:
\begin{align*}
    &\vdash \Code{let\;y\;=f_1\;RandomInt in e} : \tau \implies y:\Int, \Code{emp} \vdash e : \tau
    \\&\vdash \Code{let\;y\;=f_2\;RandomInt in e} : \tau \implies y\mapsto\Int \sep \Code{emp} \vdash e : \tau
\end{align*}

It is similar to $\{x \mapsto -\} \text{ free(x) } \{\text{emp}\}$ in SL, where we drop one write permission after apply $f_1$.

There are many possible syntactic choices we can make here; in this presentation we mostly consider under-return, thus $x{:}\Int \wand \Int$ by default means under-return.

\end{frame}

\begin{frame}[fragile]
    \frametitle{Function Typing}   

    \begin{figure}[!ht]
        \mprooftr{50mm}{
         $\Gamma \sep \Code{Over}(\Gamma), x{:}\tau_x  \vdash : e : \tau$ 
        }{TFunO}{
         $\Gamma \vdash \zzlam{x}{e} : x{:}\tau_x \sarr \tau$}
    \quad
        \mprooftr{50mm}{
         $\Gamma \sep \Code{Over}(\Gamma), x \mapsto \tau_x  \vdash : e : \tau$ 
        }{TFunU}{
         $\Gamma \vdash \zzlam{x}{e} : x{:}\tau_x \wand \tau$}
    \end{figure}

Here $\Code{Over}(\Gamma)$ means we keep all variables in $\Gamma$ with read permission, only make the $\tau_x$ be closed. I will omit this part for simplicity.
\end{frame}

\begin{frame}[fragile]
    \frametitle{Function Typing}   

    \begin{figure}[!ht]
    \mprooftr{40mm}{$ $
           }{TVarO}{
            $x:\tau_x \vdash x : \tau$}
    \quad
    \mprooftr{40mm}{$ $
           }{TVarU}{
            $x \mapsto \tau_x \vdash x \mapsto \tau$}
    \\[.8em]
    \mprooftr{70mm}{
            $\Delta_1 \vdash : v_f : y{:}\tau_y \sarr \tau_x$ \quad
            $\Gamma_2 \vdash : v : \tau_y$ \quad
            $\Gamma_3 \sep x \mapsto \tau_x \vdash : e : \tau$
           }{TAppO}{
            $(\Delta_1, \Gamma_2) \sep \Gamma_3 \vdash \zzlam{x}{v_f\;v}{e} : \tau$}
    \quad
    \mprooftr{70mm}{
            $\Gamma_1 \vdash : v_f : y{:}\tau_y \wand \tau_x$ \quad
            $\Gamma_2 \vdash : v \wand \tau_y$ \quad
            $\Gamma_3 \sep x \mapsto \tau_x \vdash : e : \tau$
           }{TAppU}{
            $(\Gamma_1 \sep \Gamma_2) \sep \Gamma_3 \vdash \zzlam{x}{v_f\;v}{e} : \tau$}
    \end{figure}
\end{frame}

\begin{frame}[fragile]
    \frametitle{Example}

    Assume type check the following program under type context $x\mapsto\Int$ against the type $\Int$.

    \begin{minted}{ocaml}
    P:
        let y1 = f x in
        let y2 = g x in
        y2
    \end{minted}
    
    Case 1: when
    \begin{align*}
        f{:}\Int \wand \Int, g{:}\Int \wand \Int
    \end{align*}

    According to \textsc{TFunU}, the argument should be type checked from a heap ($\Gamma_2$) disjoint from other heaps. The two arguments are both $x$, it means the type context should contain $(\Delta_1, x\mapsto\Int) \sep (\Delta_2, x\mapsto\Int)$ which is impossible.

\end{frame}

\begin{frame}[fragile]
    \frametitle{Example}
    
    Case 2: when
    \begin{align*}
        f{:}\Int \sarr \Int, g{:}\Int \sarr \Int
    \end{align*}
    where $f$ and $g$ work like two random generators.

    According to \textsc{TFunO}, we can use a single empty heap $(\Delta_1, \Gamma_2)$ to provide the type for the application.

    \begin{align*}
        &f{:}..., g{:}..., x \mapsto \Int \vdash P \mapsto \Int \iff 
        \\&(f{:}..., x: \Int) \sep f{:}..., g{:}..., x \mapsto \Int \vdash P \mapsto \Int
    \end{align*}
    Then according to \textsc{TAppO}, we just need to prove
    \begin{align*}
        &(f{:}..., g{:}..., x \mapsto \Int) \sep y_1\mapsto\Int \vdash ... \mapsto \Int
        \\&(f{:}..., g{:}..., x \mapsto \Int) \sep y_1\mapsto\Int, y_2\mapsto\Int \vdash y_2 \mapsto \Int
    \end{align*}
\end{frame}

\begin{frame}[fragile]
    \frametitle{Example}
    
    Case 3: when
    \begin{align*}
        f{:}\Int \sarr \Int, g{:}\Int \wand \Int
    \end{align*}

    Again, according to \textsc{TAppO}, we just need to prove
    \begin{align*}
        &(f{:}..., g{:}..., x \mapsto \Int) \sep y\mapsto\Int \vdash ... \mapsto \Int
    \end{align*}
    We still use the structural rules:
    \begin{align*}
        &(g{:}...) \sep (f{:}..., g{:}..., x \mapsto \Int) \sep y_1\mapsto\Int \vdash ... \mapsto \Int
    \end{align*}
    Then according to \textsc{TAppU}, we have
    \begin{align*}
        &y_1\mapsto\Int \sep y_2 \mapsto \Int \vdash y_2 \mapsto \Int
    \end{align*}
\end{frame}

\begin{frame}[fragile]
    \frametitle{Example}
    
    Case 4: when
    \begin{align*}
        f{:}\Int \sarr \Int, g{:}\Int \wand \Int
    \end{align*}

    According to the structural rules (especially the \textsc{ContractionEmp} rule):
    \begin{align*}
        &(f{:}...) \sep (f{:}..., g{:}..., x \mapsto \Int) \sep (g{:}..., x:\Int) \vdash P \mapsto \Int
    \end{align*}
    According to the \textsc{TAppU} rule, we have
    \begin{align*}
        &(g{:}..., x:\Int) \sep y_1\mapsto\Int \vdash ... \mapsto \Int
    \end{align*}
    Then according to \textsc{TAppO}, we have
    \begin{align*}
        &y_1\mapsto\Int \sep y_2 \mapsto \Int \vdash y_2 \mapsto \Int
    \end{align*}
\end{frame}

\begin{frame}[fragile]
    \frametitle{Example}
    
    Case 5: when
    \begin{align*}
        f{:}\Int \sarr \Int, g{:}\Int \wand (\Int, \Code{emp})
    \end{align*}

    According to the structural rules (especially the \textsc{ContractionEmp} rule):
    \begin{align*}
        &(f{:}...) \sep (f{:}..., g{:}..., x \mapsto \Int) \sep (g{:}..., x:\Int) \vdash P \mapsto \Int
    \end{align*}
    According to the \textsc{TAppU} rule, we have
    \begin{align*}
        &(g{:}..., x:\Int) \sep y_1\mapsto\Int \vdash ... \mapsto \Int
    \end{align*}
    Then according to \textsc{TAppO}, we have
    \begin{align*}
        & y_2 \mapsto \Int , y_1:\Int \vdash y_2 \mapsto \Int
    \end{align*}

    The type check will fail.
\end{frame}

\begin{frame}[fragile]
    \frametitle{Additional Rules}   

    \begin{figure}[!ht]
    \mprooftr{70mm}{
        $x{:}\rt{b}{\phi}, \Gamma \vdash e : \rt{b}{\vnu = x}  $
           }{Under}{
            $\Gamma \vdash e \mapsto \rt{b}{\phi} $}
    \end{figure}
\end{frame}

\begin{frame}[fragile]
    \frametitle{Track data flow}   

Now I want to explore a little bit more than my current setting, that is what can we gain from the data flow tracking?

\begin{minted}{ocaml}
fun x -> 
  let y = f x in
  let z = g x in
  (y, z)
\end{minted}

Assume that we really need to consume $x$ twice, where
\begin{align*}
    &f{:}\Int\wand\Int, g{:}\Int\wand\Int, x \mapsto \Int \vdash ... \\
    &f{:}\Int\wand\Int, g{:}\Int\wand\Int, x\mapsto\Int, y\mapsto\Int, z\mapsto\Int \vdash (y, z) \\
\end{align*}
where $x, y, z$ are entangled in an unknown way. Here we probably cannot get any reasonable result for the output pair.

\end{frame}

\begin{frame}[fragile]
    \frametitle{Track data flow}   

Consider another example:

\begin{minted}{ocaml}
fun x -> 
  let y = f x in
  let z = f x in
  (y, z)
\end{minted}

Here the two applications are the same, if $f$ is deterministic, we can say that $y = z$, which is potentially useful information for the output pair.

With more information about the functions (e.g., $f$ is deterministic), we can get more precise constraints for the entangled variables.

\end{frame}
